\section{The disclosed estimated value, and prior work}\label{sec:prior}
% Objective: What the number is and where it comes from: SEC staff guidance, not a rule. Position against Bauer et al., who verified such values on German certificates and left autocallables to future research.

\subsection{What the issuers say the number is}

The filings are explicit that the estimated value is a model output and whose
model it is. One of the supplements priced here states that the issue price
``exceeds the estimated value of your notes as of the time the terms of your
notes are set on the trade date, as determined by reference to GS\&Co.'s
pricing models and taking into account our credit spreads,'' and that those
``pricing models are proprietary and rely in part on certain assumptions about
future events, which may prove to be incorrect.''\footnote{Quotations in this
subsection are from the pricing supplement for note N4, filed by GS Finance
Corp.\ under accession number 0001193125-26-419171 and preserved in full in the
project repository; the filings for all six notes are identified in
Appendix~\ref{app:terms}. The other three issuers in the sample describe the
figure in materially similar terms, each naming its own proprietary models:
Morgan Stanley's supplements say the estimated value ``is determined using our
own pricing and valuation models'' and that these ``are proprietary,'' and
BofA's say it is ``determined as of the pricing date by reference to our and
our affiliates' pricing models.''} The same document lists the inputs: ``our credit spreads, interest rates (forecasted, current and
historical rates), volatility, price-sensitivity analysis and the time to
maturity of the notes.''

That list is worth reading closely, because correlation does not appear in it.
The note it describes pays on the weakest of three equity indices, and
\S\ref{sec:results} shows the investor is long correlation on every note in the
sample. Morgan Stanley's supplements list their inputs too, naming ``market inputs and
assumptions relating to the underliers, instruments based on the underliers,
volatility and other factors including current and expected interest rates,''
and correlation is absent there as well. Both lists are introduced as partial,
with ``including principally'' and ``other factors,'' so neither purports to be
complete and no inference follows that either model ignores correlation. What does follow is narrower and is the point of this paper: the
disclosure names the inputs it considers principal, and the input that most
distinguishes a worst-of structure from a single-underlying one is not among
them.

The gap between the issue price and the estimated value is itself attributed,
in the same document, to ``the underwriting discount and commissions, the
expenses incurred in creating, documenting and marketing the notes, and an
estimate of the difference between the amounts we pay to GS\&Co.\ and the
amounts GS\&Co.\ pays to us in connection with your notes.'' The first two
components appear on the cover as numbers. The third does not. It is the
internal transfer between the issuing entity and its distribution affiliate,
and it is the component that the phrase ``internal funding rate'' is usually
reaching for. That phrase belongs to the issuers rather than to the guidance. The 2012
comment letter discussed below does not contain it, while JPMorgan's supplements
in this sample carry a risk factor headed ``Estimated Value of the Notes Is
Derived by Reference to an Internal Funding Rate.'' The attribution matters
because the rate it names is the one input a reader cannot recover.

One further sentence in the same risk factor deserves attention, since it
anticipates this paper's exercise and declines to bound it. The actual value an
investor could realise ``may differ, perhaps materially, from the estimated
value of your notes determined by reference to our models due to, among other
things, any differences in pricing models or assumptions used by others.'' The
filing therefore tells the reader that another party's model may disagree, by
an amount it does not characterise, for reasons it does not enumerate. Putting
a number on that disagreement for a specific class of notes, from inputs the
reader can obtain, is what follows in \S\ref{sec:results}.

\subsection{Where the disclosure comes from}

The practice originates in an April 2012 sample comment letter from the
Division of Corporation Finance, addressed to financial institutions issuing
structured notes \citep{SEC2012}. On the valuation figure the letter asks that
issuers ``should consider prominently disclosing'' it, language that is
advisory in form as well as in effect. A sample comment letter is staff
guidance: it conveys the views of the staff, not a rule adopted by the
Commission, and it binds no one. No rule has since replaced it for this figure.

The practical consequence is that there is no prescribed form, and the sample
examined here shows the range that permits. Among the final worst-of
contingent-coupon supplements filed in the sample week, five carry no point
estimate in the filed document. Three state the figure as a bound rather than a
number, with the cover reading ``Our estimated value of the notes on trade
date: Not less than the face amount of such notes.'' The remaining two carry
only the risk-factor heading and locate the figure in an accompanying index
supplement that is not part of the document filed under the
cover.\footnote{These five supplements are not among the six priced here and
are therefore not preserved in the project repository, which holds only the
documents the valuations depend on. They remain retrievable from EDGAR under
accession numbers 0001193125-26-415466, 0001193125-26-413990 and
0001193125-26-413989 for the three quoting a bound, and 0000950103-26-015173
and 0000950103-26-015172 for the two deferring to a supplement.} All five
come from one issuer. None of this is a deficiency, since nothing requires a
point estimate; it is what discretionary disclosure looks like when it is
exercised.

The timing is constrained even where the content is not. Rule 424(b)(2)
requires the priced supplement to be filed no later than the second business
day following the earlier of the date of determination of the offering price or
the date it is first used \citep{SEC424b2}. A verification exercise therefore
has a short and known window in which the terms are fixed and public, which is
what makes a single week a coherent sample rather than an arbitrary one.

\subsection{Prior work}

The literature on retail structured products has measured issuer margins for
two decades, mostly in Europe and mostly by replication. \citet{Wilkens2003}
document premia over replication cost and a life-cycle pattern in German
products; \citet{Stoimenov2005} extend the fair-pricing analysis across
equity-linked retail instruments; \citet{Entrop2016} quantify overpricing of
discount certificates and attribute a substantial share of it to issuer credit
risk; and \citet{Ruf2011} relates pricing power to credit quality. In the US,
\citet{Henderson2011} is the canonical study, establishing that these notes are
issued above model value, and \citet{Deng2011} give an early valuation
treatment of autocallable structures alongside evidence of their growth. The
methods for valuing the structures themselves are well developed:
\citet{Guillaume2015} derives analytical results covering memory coupons,
partial protection and worst-of contingencies, and \citet{Kim2019} construct
static replications of autocallable and barrier features as an alternative to
simulation.

This paper differs from all of these in its object. They measure the distance
between a price and a model value. The question here is the distance between
two model values, one of which the issuer published. The nearest precedent for
treating a disclosed figure as the thing under examination is
\citet{Jorgensen2011}, who compare overpricing against costs the issuer
disclosed at issuance and find roughly half of it unexplained. On the US
estimated-value disclosure specifically, \citet{Lai2022} assembles a large
EDGAR sample of auto-callable reverse convertibles, but studies realised
returns rather than whether the disclosed figure can be reproduced.

The closest published work is \citet{Bauer2020}, who check issuer-disclosed
fair values against standard option pricing models for German retail
certificates and conclude that the disclosed values can be reproduced. Their
sample is of simpler instruments, and they say so: ``[m]ore complex products
like autocallables are left for future research.'' Three features of the setting
examined here make that an unusually clean gap to occupy. The product is a
worst-of on several correlated underlyings rather than a single one. On five of
the six notes redemption before maturity is not automatic, three being at the
issuer's option and two following the disclosed rational-exercise rule quoted
in \S\ref{sec:intro}, so valuation requires an optimal-stopping problem solved
against the investor rather than a barrier check. And the figure rests on staff
guidance that prescribes neither a method nor a form, which is what makes
reproducing it a question about the disclosure and not only about the model.
Whether a disclosed value can be reproduced is therefore open in this market
rather than settled elsewhere and carried across.
